% This is text associated with the open textbook "Open Electromagnetics"
% See immediately below for author, date
% <License info goes here (boilerplate, not read by project compiler)>

%ORB TITLE Electric Flux Density  
%ORB AUTHOR 0       # S.W. Ellingson, Virginia Tech (ellingson@vt.edu)
%ORB DATE   20180731 
%ORB ABSTRACT Introduction to electric flux density

\label{m0011_Electric_Flux_Density} % <-- Must match name of this file and module directory name.  Don't mess with this!
\index{electric flux density|(}
\index{field!electric|(}
\index{flux density, electric|see{electric flux density}}

Electric flux density, assigned the symbol ${\bf D}$, is an alternative to electric field intensity (${\bf E}$) as a way to quantify an electric field.
This alternative description offers some actionable insight, as we shall point out at the end of this section.

First, what is electric flux density?
Recall that a particle having charge $q$ gives rise to the electric field intensity
%
\begin{equation}
{\bf E} = \hat{\bf R} ~ q ~ \frac{1}{4\pi R^2} ~ \frac{1}{\epsilon}
\label{m0011_eEparticle1}
\end{equation}
%
where $R$ is distance from the charge and $\hat{\bf R}$ points away from the charge.
Note that ${\bf E}$ is inversely proportional to $4\pi R^2$, indicating that ${\bf E}$ decreases in proportion to the area of a sphere surrounding the charge.
Now integrate both sides of Equation~\ref{m0011_eEparticle1} over a sphere $\mathcal{S}$ of radius $R$:
%
\begin{equation}
\oint_{\mathcal{S}} {\bf E}({\bf r}) \cdot d{\bf s} = \oint_{\mathcal{S}} \left[ \hat{\bf R} ~ q ~ \frac{1}{4\pi R^2} ~ \frac{1}{\epsilon} \right] \cdot d{\bf s}
\end{equation}
%
Factoring out constants that do not vary with the variables of integration, the right-hand side becomes:
%
\begin{equation}
q ~ \frac{1}{4\pi R^2} ~ \frac{1}{\epsilon} ~ \oint_{\mathcal{S}} \hat{\bf R} \cdot d{\bf s}
\nonumber
\end{equation}
%
Note that $d{\bf s}=\hat{\bf R}ds$ in this case, and also that $\hat{\bf R}\cdot\hat{\bf R}=1$.  Thus, the right-hand side simplifies to:
%
\begin{equation}
q ~ \frac{1}{4\pi R^2} ~ \frac{1}{\epsilon} ~ \oint_{\mathcal{S}} ~ds
\nonumber
\end{equation}
%
The remaining integral is simply the area of $\mathcal{S}$, which is $4\pi R^2$.  Thus, we find:
%
\begin{equation}
\oint_{\mathcal{S}} {\bf E}({\bf r}) \cdot d{\bf s} = \frac{q}{\epsilon}
\label{m0011_eEparticle2}
\end{equation}
%
The integral of a vector field over a specified surface is known as \emph{flux} 
\index{flux}
(see ``Additional Reading'' at the end of this section). 
Thus, we have found that the flux of ${\bf E}$ through the sphere $\mathcal{S}$ is equal to a constant, namely $q/\epsilon$.
Furthermore, this constant is the same regardless of the radius $R$ of the sphere.
Said differently, the flux of ${\bf E}$ is constant with distance, and does not vary as ${\bf E}$ itself does.  
Let us capitalize on this observation by making the following small modification to Equation~\ref{m0011_eEparticle2}:
%
\begin{equation}
\oint_{\mathcal{S}} \left[\epsilon {\bf E}\right] \cdot d{\bf s} = q
\end{equation}
%
In other words, the flux of the quantity $\epsilon {\bf E}$ is equal to the enclosed charge, regardless of the radius of the sphere over which we are doing the calculation.
This leads to the following definition:

%ORB HIGHLIGHT
\begin{mdframed}[backgroundcolor=yellow!20,nobreak=true] 
The electric flux density 
\index{electric flux density!definition}
${\bf D} = \epsilon {\bf E}$, having units of C/m$^2$, is a description of the electric field in terms of flux, as opposed to force or change in electric potential. 
\end{mdframed}

It may appear that ${\bf D}$ is redundant information given ${\bf E}$ and $\epsilon$, but this is true only in homogeneous media.
\index{homogeneous (media)}
The concept of electric flux density becomes important -- and decidedly not redundant -- when we encounter boundaries between media having different permittivities.  
As we shall see in Section~\ref{m0021_Boundary_Conditions_on_D}, boundary conditions on ${\bf D}$ constrain the component of the electric field that is perpendicular to the boundary separating two regions.  If one ignores the ``flux'' character of the electric field represented by ${\bf D}$ and instead considers only ${\bf E}$, then only the \emph{tangential} component of the electric field is constrained.  
In fact, when one of the two materials comprising the boundary between two material regions is a perfect conductor, 
\index{conductor!perfect}
\index{perfect conductor|see{conductor, perfect}}
then the electric field is \emph{completely determined} by the boundary condition on ${\bf D}$.  This greatly simplifies the problem of finding the electric field in a region bounded or partially bounded by materials that can be modeled as perfect conductors, including many metals.  In particular, this principle makes it easy to analyze capacitors. 
 
We conclude this section with a warning.
Even though the SI units for ${\bf D}$ are C/m$^2$, ${\bf D}$ describes an electric field and \emph{not} a surface charge density.
It is certainly true that one may describe the amount of charge distributed over a surface using units of C/m$^2$.
However, ${\bf D}$ is not necessarily a description of \emph{actual} charge, and there is no implication that the source of the electric field is a distribution of surface charge.
On the other hand, it \emph{is} true that ${\bf D}$ can be interpreted as an \emph{equivalent} surface charge density that would give rise to the observed electric field, and in some cases, this equivalent charge density turns out to be the actual charge density. 

%%%%%%%%%%%%%%%%%%%%%%%%%%

{\bf Additional Reading:}
\begin{itemize}
\item ``\href{https://en.wikipedia.org/wiki/Flux}{Flux}'' on Wikipedia.
\end{itemize}

\index{electric flux density|)}
\index{field!electric|)}
